\chapter{Framework for Pedestrian Detection}

\section{Feature Extraction}
Notes may be presented as footnotes or chapter endnotes according to the common practice of the discipline. Sources, permissions, and explanatory notes should be cited clearly and consistently.

\subsection{HOG Features}
Each table should have a table number and title. The caption should be placed above the table, and the table should appear after it is mentioned in the text. \Cref{tab:video-info} gives a sample table.

\begin{table}[htbp]
  \centering
  \caption{Effect information of detection videos}
  \label{tab:video-info}
  \vspace{0.2cm}
  \begin{tabular}{lccc}
    \toprule
    Video Sequence & Frame Size & Avg. Detection Time & Avg. Accuracy \\
    \midrule
    Video 1 & 640 x 480 & 62 ms/frame & 93.3\% \\
    Video 2 & 640 x 480 & 65 ms/frame & 92.7\% \\
    Video 3 & 320 x 240 & 44 ms/frame & 94.1\% \\
    \bottomrule
  \end{tabular}
\end{table}
